% ============================================================
% Chapter: Conclusion and Future Work
% Thesis: Solving the FSRCPSP Using Hybrid ML Approaches
% Author: Souvik Chakraborty
% ============================================================
\chapter{Conclusion and Future Work}
\label{chap:conclusion}

% ============================================================
\section{Summary of Contributions}
\label{sec:conclusion_summary}
% ============================================================
This thesis addressed the Flexible Stochastic
Resource-Constrained Project Scheduling Problem (FSRCPSP),
a recently introduced extension of the classical RCPSP that
combines flexible resource profiles with stochastic workloads.
The SAA-FSRCPSP model grows as
$O(|V| \cdot |K| \cdot |T| \cdot |S|)$, making exact solution
intractable beyond small instances and motivating the
development of decomposition heuristics. The thesis made four
contributions, each corresponding to one of the research
questions posed in Chapter~\ref{chap:introduction}.

\textbf{Contribution~1: Group-based activity-oriented
Relax-and-Fix.}
Chapter~\ref{chap:grouprf} introduced a group-based R\&F
heuristic that replaces the single-activity decomposition of
the baseline (Chapter~\ref{chap:baseline}) with an unsupervised
machine learning pipeline. Ward hierarchical clustering
partitions activities into precedence-safe groups based on
extracted features, and each group is fixed simultaneously in a
single MIP subproblem. The experimental evaluation
(Chapter~\ref{chap:experiments}) showed that this grouping
yields consistent speedups of $3$--$12\times$ on j10 instances
and $4$--$5\times$ on j30 instances, while matching or improving
solution quality. At $|K|=4$ on j10, the group-based method
recovered feasibility on all 7 instances the baseline missed
($p = 0.016$, McNemar test), and at $|K|=1$ on j30 it solved
4 additional instances that the baseline could not reach.

\textbf{Contribution~2: Feature pool and ablation study.}
A candidate pool of $18 + 8|K|$ features ($50$ for $|K|=4$)
was designed, spanning network topology (Groups~A--B), duration
(Group~C), scheduling urgency (Group~D), timing (Group~E), and
stochastic resource load statistics (Groups~F--G). A systematic
ablation over 32 configurations on j60 instances with $|K|=4$
(Section~\ref{sec:ablation}) revealed that the four network
position features (precedence level, in-degree, out-degree,
longest path to sink) are Pareto-optimal: no alternative
configuration achieved both a lower gap and equal or higher
feasibility. Adding resource statistics to the feature set
could reduce the mean gap from $8.09\%$ to $7.72\%$, but at
the cost of losing feasible instances and increasing runtime
by up to $2.9\times$.

\textbf{Contribution~3: Scalability extensions.}
Two extensions were developed to push the method beyond j30
instances. Time aggregation
(Section~\ref{sec:alpha_sensitivity}) compresses the time
horizon from $|T|$ to $\lceil |T|/\alpha \rceil$ periods,
reducing per-subproblem model size and enabling 100\%
feasibility on j60 instances with $\alpha \in \{5, 7\}$.
Per-group scenario reduction
(Section~\ref{sec:scen_red_exp}) applies PAM K-medoids
separately to each group's workload vectors, selecting the
$n_{\mathrm{keep}} = 3$ most representative scenarios per
subproblem. On j120 instances with $|K|=4$, per-group
reduction achieved a mean gap of $10.30\%$ versus $31.83\%$
for the best global method, a $3.1\times$ improvement, because
each group's subproblem uses scenarios tailored to its own
resource profile. Across all $|K|$ values, per-group reduction
produced \emph{negative} mean gaps ($-13\%$ to $-17\%$),
meaning the MIP solver consistently improved upon the SSGS
warm-start.

\textbf{Contribution~4: Systematic evaluation.}
The evaluation covered PSPLIB instances with
$n \in \{10, 30, 60, 120\}$ activities and
$|K| \in \{1, 2, 3, 4\}$ resource types, totalling over
1\,600 instance-method combinations. Statistical validation
via Wilcoxon signed-rank tests confirmed significant runtime
improvements in five of eight j10/j30 configurations, and
Dolan--Mor\'e performance profiles showed that the group-based
method is the faster solver on over $75\%$ of the 323 jointly
feasible instances.

% ============================================================
\section{Answers to the Research Questions}
\label{sec:conclusion_rq}
% ============================================================
The four research questions from
Section~\ref{sec:problem_statement} can now be answered.

\textbf{RQ\,1:} Can unsupervised machine learning guide the
R\&F decomposition? Yes. Ward hierarchical clustering on
network position features produces precedence-safe groups that
give the MIP solver richer combinatorial context per
subproblem. The resulting grouping is stable across instance
sizes and resource counts, typically collapsing j10 instances
into a single group and partitioning j30 instances into
2--3 groups.

\textbf{RQ\,2:} Does grouping improve over the baseline?
Yes, in both runtime and feasibility. Speedups range from
$3\times$ to $12\times$ on j10 and $4$--$5\times$ on j30.
Feasibility is equal or better in seven of the eight
j10/j30 configurations tested, with statistically significant
improvement at j10/$|K|=4$. Solution quality is effectively
identical: the baseline wins on fewer than $1\%$ of jointly
feasible instances.

\textbf{RQ\,3:} Which features matter most?
The four network position features (Group~A) are sufficient
and Pareto-optimal. Resource statistics can marginally improve
solution quality but at the cost of feasibility and runtime.
Correlation-based pruning ($\theta_{\mathrm{corr}} = 0.95$)
ensures that the full 50-feature set degrades gracefully to
match the net4 baseline, but provides no additional benefit.

\textbf{RQ\,4:} Do the scalability extensions work?
Yes. Time aggregation with $\alpha = 2$ raises j30 feasibility
from $52$--$78\%$ to $92$--$100\%$. For j60, $\alpha \in \{5,7\}$
achieves 100\% feasibility across all $|K|$. Per-group scenario
reduction is essential for j120: it reduces the mean gap from
$16$--$23\%$ (baseline) to $-13\%$ to $-17\%$ (negative,
improving on the SSGS warm-start) and raises feasibility at
$|K| \geq 3$.

% ============================================================
\section{Limitations}
\label{sec:conclusion_limitations}
% ============================================================
Several limitations should be noted. All stochastic workloads
are generated from a single $\mathrm{Beta}(0.5,\,2.25)$
distribution; results may differ under heavier-tailed or
correlated distributions. The SSGS warm-start contains a known
inconsistency in worst-case scenario selection
(Chapter~\ref{chap:baseline}), and correcting it could shift
the relative performance of both methods. The j60 and j120
results depend on time aggregation and scenario reduction,
respectively, so the reported gaps mix algorithmic and
approximation effects. All runtimes are wall-clock measurements
on a single Apple~M1 machine running Gurobi under Rosetta~2
emulation and are not directly comparable to native x86\_64
benchmarks. Finally, the ablation was conducted on 10 j60
instances with $|K|=4$; while the results are consistent with
the full 50-instance experiments, a larger ablation sample
would strengthen the conclusions.

% ============================================================
\section{Future Work}
\label{sec:future_work}
% ============================================================
Several directions emerge for extending the work presented in
this thesis.

\paragraph{Adaptive and instance-specific grouping.}
The current pipeline applies a fixed feature subset and a
single clustering to each instance. An adaptive scheme could
select features or adjust the number of groups based on
instance characteristics detected at runtime, for example by
training a meta-learner on the ablation data to predict which
feature configuration works best for a given instance structure.

\paragraph{Alternative clustering methods.}
Ward hierarchical clustering was chosen for its determinism and
interpretability, but other methods may be worth exploring.
Spectral clustering on the precedence graph Laplacian would
directly exploit the network structure, and density-based
methods such as HDBSCAN could identify groups of varying size
without requiring a fixed distance threshold.

\paragraph{Correlated and heavy-tailed scenario generation.}
The current $\mathrm{Beta}(0.5,\,2.25)$ workload model assumes
independence across activities and scenarios. In practice,
project risks are often correlated (e.g.\ weather affects all
outdoor activities simultaneously). Evaluating the method under
copula-based or factor-model scenario generators would test its
robustness to more realistic uncertainty structures.

\paragraph{Multi-objective formulation.}
The current objective minimises expected makespan. In many
applications, project managers must balance makespan against
resource cost, risk (e.g.\ CVaR of the makespan distribution),
or schedule robustness. Extending the SAA-FSRCPSP to a
multi-objective formulation and adapting the R\&F decomposition
to handle Pareto-optimal trade-offs is a natural next step.

\paragraph{Improved warm-start heuristics.}
The SSGS warm-start uses a single priority rule and a
worst-case scenario selection that is dimensionally
inconsistent across resources
(Chapter~\ref{chap:baseline}). Replacing it with a
multi-start heuristic, a genetic algorithm, or a reinforcement
learning policy could improve the initial solution quality and
thereby reduce the work left for the MIP solver.

\paragraph{Interaction of time aggregation and scenario
reduction.}
The j120 experiments apply both time aggregation ($\alpha=15$)
and per-group scenario reduction ($n_{\mathrm{keep}}=3$)
simultaneously, but the interaction between these two
approximations has not been characterised. A factorial
experiment varying both parameters jointly would clarify
whether the two techniques are complementary or partially
redundant, and could identify configurations that achieve
better gap-feasibility trade-offs.

\paragraph{Real-world instances.}
All experiments use PSPLIB-generated instances with synthetic
stochastic workloads. Validating the approach on real-world
project data, for example from construction, pharmaceutical
development, or software engineering, would test whether the
structural assumptions (precedence-based grouping, Beta-distributed
workloads, fixed resource capacities) hold in practice.

\paragraph{Exact solution methods for small instances.}
For j10 instances, the group-based method already solves the
full MIP in a single pass. For slightly larger instances
($n = 15$--$20$), branch-and-cut or branch-and-price methods
might close the gap to optimality, providing reference solutions
against which heuristic quality can be measured more precisely
than relative to the SSGS warm-start.
